\subsection{卡西米尔元}

命题 11. 设 \(\mathfrak{g}\) 是域 K 上的李代数，U 是其包络代数，\(\mathfrak{h}\) 是 \(\mathfrak{g}\) 的有限维理想，\(\beta\) 是 \(\mathfrak{g}\) 上的不变双线性形式，并且其在 \(\mathfrak{h}\) 上的限制非退化。设 \((e_i)_{1\leqslant i\leqslant n}, (e_j')_{1\leqslant j\leqslant n}\) 是 \(\mathfrak{h}\) 的两个基，满足 \(\beta(e_i,e_j') = \delta_{ij}\)。

则 U 中的元素 \(c = \sum_{i=1} e_i e_i'\) 属于 U 的中心，并且与基 \((e_i)\) 的选择无关。

对 \(x \in \mathfrak{g}\)，令 \(x_{\mathfrak{h}}\) 为 \(\mathfrak{h}\) 在 \(\operatorname{ad}_{\mathfrak{g}} x\) 上的限制。则 \(x \mapsto x_{\mathfrak{h}}\) 是 \(\mathfrak{g}\) 在向量空间 \(\mathfrak{h}\) 上的表示，而 \(\beta'\) 在 \(\beta\) 上的限制 \(\mathfrak{h}\) 在此

\[
\sum_ {i = 1} ^ {n} e _ {i} \otimes e _ {i} ^ {\prime}
\]

 表示。根据第 5 号例 3，张量 \(\sum_{i=1}^{n} e_i \otimes e_i'\) 与基 \((e_i)\) 的选择无关，并且是 \(\mathfrak{h}\) 的张量代数中的不变元。它也是 \(T\) 的张量代数 \(\mathfrak{g}\) 中的元素；对于由 \(\mathfrak{g}\) 的伴随表示导出的表示，它是不变的。因此它在 U 中的典范像，即 \(c\)，与基 \((e_i)\) 的选择无关，并且对于第 2 号末尾所考虑的 \(\mathfrak{g}\) 在 U 上的表示，它是不变元。因此，该元素与 \(\mathfrak{g}\) 的每个元素可交换，从而属于 U 的中心。

当 \(\beta\) 是与 g-模 M 关联的双线性形式时，命题 11 中的元素 \(c\) 称为与 M（或与相应表示）关联的卡西米尔元。当 \(\beta\) 在 \(\mathfrak{h}\) 上的限制非退化时，此元素存在。

命题 12。设 g 是域 K 上的李代数，h 是 \(n\) 维理想，\(\mathbf{M}\) 是在 \(\mathbf{K}\) 上有限维的 g-模。设 \(c\) 是与 \(\mathbf{M}\) 和 \(\mathfrak{h}\) 关联的卡西米尔元（假定其存在）。

\begin{enumerate}
\def\labelenumi{(\alph{enumi})}
\item
  \(\mathrm{Tr}(c_{\mathrm{M}})=n.\)
\item
  若 \(\mathbf{M}\) 是简单模，且 \(n\) 不被 \(\mathbf{K}\) 的特征整除，则 \(c_{\mathbf{M}}\) 是 \(\mathbf{M}\) 的自同构。
\end{enumerate}

采用命题 11 的记号，

\[
\mathrm{Tr} (c _ {\mathrm{M}}) = \sum_ {i = 1} ^ {n} \mathrm{Tr} ((e _ {i}) _ {\mathrm{M}} (e _ {i} ^ {\prime}) _ {\mathrm{M}}) = \sum_ {i = 1} ^ {n} \beta (e _ {i}, e _ {i} ^ {\prime}) = n.
\]

因此，若 \(n\) 不被 \(\mathbf{K}\) 的特征整除，则 \(c_{\mathbf{M}} \neq 0\)。另一方面，由于 \(c\) 属于 \(\mathrm{U}\) 的中心，\(c_{\mathbf{M}}\) 与所有 \(x_{\mathbf{M}}\)（\(x \in \mathfrak{g}\)）可交换。若进一步假设 \(\mathbf{M}\) 简单，则 \(c_{\mathbf{M}}\) 在 \(\mathcal{L}(\mathbf{M})\) 中可逆（《代数》，第 VIII 章，§ 4，第 3 号命题 2）。

